
\documentclass[a4paper,12pt]{letter}
\usepackage[utf8]{inputenc}
\usepackage[english,slovene]{babel}
%\usepackage{blindtext}
%\usepackage{pdflscape}
\usepackage[tikz]{bclogo}
\usepackage{qrcode,marvosym}
\usepackage{pgf,tikz,pgfplots}
\pgfplotsset{compat=1.14}
\usepackage{mathrsfs}
\usetikzlibrary{arrows}
\usepackage{graphicx}
\newcommand{\ds}{\displaystyle}
\usepackage{fancybox}
\usepackage{amsfonts}
\usepackage{amsmath}
%\usepackage{color}
\usepackage{wallpaper}
\usepackage{hyperref}
\usepackage[cp1250]{inputenc}
\usepackage{array}%%%%%%%%%%%%%%%%%%%%%%%%%%%%%%%%%%%%%%%%%%%%
\newcolumntype{M}[1]{>{\centering\arraybackslash}m{#1}}%%%%
\newcolumntype{N}{@{}m{0pt}@{}}%%%%%%%%%%%%%%%%%%%%%%%%%%%%
   \usepackage{fancyhdr,enumerate}

 \newcommand{\ora}{\overrightarrow}
 \renewcommand{\headrulewidth}{3pt}
 \renewcommand{\headrule}{\hbox to\headwidth{%
   \color{gray}\leaders\hrule height \headrulewidth\hfill}}
    \renewcommand{\footrulewidth}{2pt}
    \renewcommand{\footrule}{\hbox to\headwidth{%
  \color{brown}\leaders\hrule height \footrulewidth\hfill}}
\usepackage{gensymb}
\usepackage{amssymb} 

  \usepackage{fancyhdr}
  \textwidth 20 true cm
  \oddsidemargin -2 true cm
  \evensidemargin -2 true cm
  \topmargin -3.5 true cm
  \textheight 27.5 true cm
  \linespread{1.2}
    \renewcommand{\headrulewidth}{0.3pt}
   \renewcommand{\footrulewidth}{1.9pt}
  \definecolor{qqwuqq}{rgb}{1,0,0 }
  \definecolor{qqqqff}{rgb}{0,0,1}
  \definecolor{qqffqq}{rgb}{0,1,0}
  \definecolor{ffqqqq}{rgb}{1,0,0}
\definecolor{zzuxux}{rgb}{0,0,0}%{0.92,0.03,0.03}
\usepackage{mdframed}
\usepackage{multicol}
 \definecolor{bgblue}{RGB}{0,0,253}
\usepackage{xcolor}
\usepackage{eurosym}
%%%%%%%%%%%%%%%%%%%%%%%%%%%%%%%%%%%%%%%%%%%%%%%%%%%%%%%%
\newcounter{tocke}
\setcounter{tocke}{0}
\usepackage{tikz-3dplot}
%%%%%%%%%%%%%%%%%%%%%%%%%%%%%%%%%%%%%%%%%%%%%%%%%%%%%%%%%%
\mdfdefinestyle{theoremstyle}{%
linecolor=brown!40,linewidth=1pt,%
frametitlerule=true,%
frametitlebackgroundcolor=brown!50,
innertopmargin=\topskip,
}

\mdfdefinestyle{theoremstyle1}{%
linecolor=brown,middlelinewidth=1pt,%
frametitlerule=true,%
apptotikzsetting={\tikzset{mdfframetitlebackground/.append style={%
shade,left color=black!40, right color=gray!10},
mdfbackground/.append style={
top color=white!100!white,
bottom color=black!5
},
}}, 
frametitlerulecolor=gray,
frametitlerulewidth=2pt,
innertopmargin=\topskip,
innerbottommargin=\topskip,
}
\mdtheorem[style=theoremstyle1]{naloga}{Naloga}
\mdtheorem[style=theoremstyle]{resitev}{Analiza Naloge}
\mdtheorem[style=theoremstyle]{kriterij}{\v{S}tevilo dose\v{z}enih to\v{c}k na testu}
%%%%%%%%%%%%%%%%%%%%%%%%%%%%%%%%%%%%%%%%%%%%%%%%%%%%%%%%%%
\usepackage[table]{xcolor}
\usepackage{titlesec}
\usepackage{venndiagram}
\usepackage[printwatermark]{xwatermark}
\usepackage{xcolor}
\newwatermark[allpages,color=brown!2,angle=45,scale=5,
xpos=-5,ypos=0]{matej.info}


%%%%%%%%%%%%%%%%%%%%%%%%%%%%%%%%%%%%%%%%%%%%%%%%%%%%%%%%%%
\titleformat{\subsection}
  {\normalfont\Large\bfseries\color{brown}}
  {\thesubsection}
  {1em}
  {}

\titleformat{\subsubsection}
  {\normalfont\large\bfseries\color{brown}}
  { \thesubsubsection}
  {1em}
  {}
%%%%%%%%%%%%%%%%%%%%%%%%%%%%%%%%%%%%%%%%%%%%%%%%%%%%%%%%%%
\begin{document}
\pagestyle{fancy}
\renewcommand\headrulewidth{0pt}
\lhead{}\chead{}\rhead{}
\rfoot{\vspace*{-2\baselineskip}}


\renewcommand{\vec}{\overrightarrow}

%%%%%%%%%%%%%%%%%%%%%%55\Large

\everymath{\displaystyle}
\begin{minipage}{20cm}
\begin{bclogo}[couleur=brown!40,
  arrondi=0,ombre=true, couleurOmbre=gray!50,
logo =\bcplume,
  noborder=true]
{\textrm{{\Large{Test 2.3} \Huge{}:
%
%
%%%%%%%%%%%%%%%%%%%%%%%%%%%%%%%%%%%%%%%%%
%%%%%%%%%%%%%%%%%%%%%%%%%%%%%%%%%%%%%%%%%
% TUKAJ VNESI VSEBINO TESTA
\large{Eksponentna in logaritemska funkcija. Zaporedja.} \hfill $PTI-2$}}}
%%%%%%%%%%%%%%%%%%%%%%%%%%%%%
\end{bclogo}
\end{minipage}
\begin{minipage}{20cm}
\begin{bclogo}[couleur=brown!40,ombre=true, couleurOmbre=gray!50,
  arrondi=0,
logo =\bcplume,
  noborder=true]
 \setlength{\parskip}{}{}
{\textsc{{\large{Ime in priimek:}
%%%%%%%%%%%%%%%%%%%%%%%%%%%%%%%%%%%%%%%%%
% TUKAJ VNESI VSEBINO TESTA
\vspace{0.2cm} \rule{15cm}{0.2mm}\hfill }}}
%%%%%%%%%%%%%%%%%%%%%%%%%%%%%%%%%%%%%%%%%
\end{bclogo}
\end{minipage}


%%%%%%%%%%%%%%%%%%%%%%%%%%%%%%%%%
%%%%%%%%%%%%%%%%%%%%%%%%%%%%%%%%%%%%%%%%%

\def\tocke2{3+3+3}
\begin{naloga}[\hfill \quad $\tocke2$
\qquad \rightsquigarrow \vert a.\qquad|b.\qquad|c.\qquad|] \addtocounter{tocke}{\tocke2}
Reši enačbo:
\vspace{-0.5cm}
\begin{multicols} {3}
   a) $\log_3(2x+1)=3$

b) $4^x=9^x$

c) $2^x=20$ 
\end{multicols}
\end{naloga}

%%%%%%%%%%%%%%%%%%%%%%%%%%%%%%%%%%%%%%%%%

\newpage
\def\tocke18{5+1}
\begin{naloga}[\hfill \quad $\tocke18$
\qquad \rightsquigarrow \vert a.\qquad\vert b.\qquad|
c.\qquad|] \addtocounter{tocke}{\tocke18}

a) Nariši graf funkcije $f(x)=3^{x}-2$, tako da najprej izračunaš začetno vrednost in ničlo ter določiš asimptoto.

b) Izračunaj vrednost točke $A(2,y)$, če leži na grafu ter jo označi.
\end{naloga}

\begin{tikzpicture}[line cap=round,line join=round,>=triangle 45,x=1.0cm,y=1.0cm]
\begin{axis}[
x=1.0cm,y=1.0cm,
axis lines=middle,
ymajorgrids=true,
xmajorgrids=true,
xmin=-6,
xmax=6,
ymin=-4,
ymax=8,
xtick={-6,-5,...,6},
ytick={-4,-3,...,8},]
\clip(-6,-4) rectangle (6,8);
\end{axis}
\end{tikzpicture}

%%%%%%%%%%%%%%%%%%%%%%%%%%%%%%%%%%%%%%%%%

\newpage
\def\tocke6{4+4}
\begin{naloga}[\hfill \quad $\tocke6$
\qquad \rightsquigarrow \vert a.\qquad| b.\qquad| c.\qquad| ] \addtocounter{tocke}{\tocke6}
V aritmetičnem zaporedju je prvi člen enak 5, peti člen pa 21. 

a) Izračunaj diferenco zaporedja in tretji člen.

b) Izračunaj vsoto prvih 12 členov tega zaporedja.
\end{naloga}

\vspace{10cm}

%%%%%%%%%%%%%%%%%%%%%%%%%%%%%%%%%%%%%%%%%

\def\tocke6{4+3}
\begin{naloga}[\hfill \quad $\tocke6$
\qquad \rightsquigarrow \vert a.\qquad| b.\qquad| c.\qquad| ] \addtocounter{tocke}{\tocke6}
a) Naj bo $\log a=2, \log b=3,\log c=4.$ Izračunaj $\log x,$ če je 
\[
x=\dfrac{ab^2}{\sqrt{c}}.
\]

b) Izračunaj:
\[
\log_2 16-\log_4 2+\ln 1+\log 100-\log 10
\]
\end{naloga}

%%%%%%%%%%%%%%%%%%%%%%%%%%%%%%%%%%%%%%%%%

\newpage
\def\tocke3{3+4+3}
\begin{naloga}[\hfill \quad $\tocke3$
\qquad \rightsquigarrow \vert a.\qquad\vert b.\qquad] \addtocounter{tocke}{\tocke3}
a) V aritmetičnem zaporedju $\rule{1cm}{0.1mm},\;4\;,9\;,\rule{1cm}{0.1mm}\;\rule{1cm}{0.1mm}, \ldots$ dopolni tri manjkajoče člene.

b) Izračunaj 8. člen in vsoto prvih 8 členov.

c) Na katerem mestu je člen $499\;?$
\end{naloga}

\vfill
\mdtheorem[style=theoremstyle1]{kriterij}{Kriterij ocenjevanja}
\begin{kriterij*}[\hfill \v{s}tevilo vseh to\v{c}k na testu: $\arabic{tocke}$]
 $$
 \begin{array}{||c|c|c|c|c|c||
>{}c|c|}
 \hline\hline
 \textup{ocena}& 1&2&3&4&5&\textrm{uspe\v{s}nost v } \%& \textrm{\bf OCENA}\\
  \hline\hline
 \% & [0,45) &[45,60)&[60,75
 ) &[75,90)    &    [90,100]& \mbox{\phantom{ \huge{$\frac{5}{445}$
 15}}
 } 
 & \\
 \hline
 \end{array} \hspace*{0.5cm}\vspace*{-0.3cm} \qrcode[hyperlink,height=0.8in]{http://www.matej.info}
$$
\vspace*{0.1cm}
\end{kriterij*}
%%%%%%%%%%%%%%%%%%%%%%%%%%%%%%%%%%%%%%%%%
\newpage
\subsection{Rešitve:}

\subsubsection{1. naloga}

a)
\[
2x+1=27 \Rightarrow x=13
\]

b)
\[
\left(\frac{4}{9}\right)^x=1 \Rightarrow x=0
\]

c)
\[
x=\log_2 20=\frac{\log 20}{\log 2}
\]

%%%%%%%%%%%%%%%%%%%%%%%%%%%%%%%%%%%%%%%%%

\subsubsection{2. naloga}

a)
\[
f(0)=1-2=-1
\]

\[
3^x-2=0 \Rightarrow 3^x=2 \Rightarrow x=\log_3 2
\]

\[
y=-2
\]

b)
\[
y=f(2)=9-2=7
\]

\[
A(2,7)
\]

%%%%%%%%%%%%%%%%%%%%%%%%%%%%%%%%%%%%%%%%%

\subsubsection{3.naloga }

\[
21=5+4d \Rightarrow d=4
\]

\[
a_3=5+2\cdot4=13
\]

\[
S_{12}=6(10+44)=324
\]

%%%%%%%%%%%%%%%%%%%%%%%%%%%%%%%%%%%%%%%%%

\subsubsection{4.naloga}

a)
\[
\log x=2+2\cdot3-\frac{1}{2}\cdot4=6
\]

b)
\[
4-\frac12+1=\frac{9}{2}
\]

%%%%%%%%%%%%%%%%%%%%%%%%%%%%%%%%%%%%%%%%%

\subsubsection{5. naloga}

\[
d=5
\]

\[
-1,\;4,\;9,\;14,\;19
\]

b)
\[
a_8=34,\quad S_8=132
\]

c)
\[
n=101
\]

%%%%%%%%%%%%%%%%%%%%%%%%%%%%%%%%%%%%%%%%%

\end{document}